% Source: https://artofproblemsolving.com/community/c3h2493040_sum_of_integers__mathcounts
% Stable identity: aops-topic-2493040-q01
\textbf{(a)} A sum of $n$ consecutive integers starting at $x$ is $n$ times the average of the first and last terms:
\[x+(x+1)+\cdots+(x+n-1)=\dfrac{n(2x+n-1)}{2}.\]
So we need $n(2x+n-1)=30030=2\cdot3\cdot5\cdot7\cdot11\cdot13$ with $n\ge2$ and $x\ge1$. Put $m=2x+n-1$; then $nm=30030$ and $m-n=2x-1$.

Conversely, take any factorisation $30030=nm$ with $2\le n<m$. Since $30030$ contains exactly one factor $2$, one of $n,m$ is even and the other odd, so $m-n$ is odd and $x=\dfrac{m-n+1}{2}$ is a positive integer. So the representations correspond exactly to factorisations $30030=nm$ with $2\le n<m$.

$30030$ has $2^6=64$ divisors; it is not a perfect square, so they form $32$ pairs $\{n,m\}$ with $n<m$. Excluding $n=1$ (a single term) leaves $\boxed{31}$ representations.

\textbf{(b)} The student's equation is wrong: the sum of the $n$ terms is the \emph{number of terms} $n$ times the average $\dfrac{x+(x+n-1)}{2}$, not the first term $x$ times the average. The count $2^6-1$ also counts every divisor of $30030$ except one, but each representation corresponds to an \emph{unordered} pair of complementary divisors $n<m$; the condition $x\ge1$ forces $m>n$, which halves the count, and then the case $n=1$ is removed: $32-1=31$.
